\documentclass[11pt,a4paper]{report}
\usepackage[T1]{fontenc}
\usepackage[utf8]{inputenc}
\usepackage[spanish,es-nodecimaldot]{babel}
\usepackage[margin=24mm]{geometry}
\usepackage{graphicx,xcolor,booktabs,tabularx,float,hyperref}
\definecolor{verde}{HTML}{184B3A}
\hypersetup{colorlinks=true,linkcolor=verde,urlcolor=verde,pdftitle={Manual de usuario - Laboratorio de Grafos y Búsquedas},pdfauthor={ticonsky}}
\graphicspath{{capturas/}}
\setlength{\parindent}{0pt}
\setlength{\parskip}{.45em}
\newcommand{\boton}[1]{\textbf{#1}}
\newcommand{\pasos}{\begin{enumerate}}
\newcommand{\finpasos}{\end{enumerate}}
\newcommand{\captura}[3]{\begin{figure}[H]\centering
\includegraphics[width=\linewidth,height=.49\textheight,keepaspectratio]{#1.png}
\caption{#2}\label{fig:#3}\end{figure}}
\newcommand{\metodo}[3]{\subsection{#1}\textbf{Datos y acción.} #2\par\textbf{Lectura del resultado.} #3}
\begin{document}
\begin{titlepage}
\color{verde}
\rule{\textwidth}{1.2pt}
\vspace*{22mm}

{\Large Ciencias de la Computación 2\par}
\vspace{12mm}
{\Huge\bfseries Manual de usuario\par}
\vspace{5mm}
{\LARGE Laboratorio de Grafos y Búsquedas\par}
\vspace{18mm}
{\large Universidad Distrital Francisco José de Caldas\par}
\vfill
{\large Desarrollo del simulador: ticonsky\par}
{\large Noviembre de 2026\par}
\vspace{5mm}
\rule{\textwidth}{1.2pt}
\end{titlepage}
\setcounter{tocdepth}{1}
{\small\tableofcontents}

\chapter{Introducción}
\section{Propósito y alcance}
El Laboratorio de Grafos y Búsquedas es un simulador para practicar los contenidos de Ciencias de la Computación 2. Permite introducir datos, aplicar algoritmos y examinar su representación gráfica o tabular. Este manual describe el uso de los 26 métodos accesibles desde los catálogos y documenta ejercicios que se pueden repetir con el botón \boton{Ejemplo}.

\section{Público y convenciones}
Está dirigido a estudiantes y docentes. Los nombres en \textbf{negrita} identifican controles de la interfaz. Los valores como \texttt{45} o \texttt{A} son datos de ejemplo. Los nombres de los controles pueden variar por módulo, especialmente cuando se pasa de búsquedas a grafos.

\chapter{Acceso y primeros pasos}
\section{Requisitos y arranque}
Para usar una copia publicada basta un navegador moderno. Para abrir una copia local por medio de un servidor, se requiere Python 3. Abre una terminal en la carpeta \texttt{SimuladorCC2}, donde está \texttt{index.html}, y ejecuta:
\begin{verbatim}
python3 -m http.server 8000 --bind 127.0.0.1
\end{verbatim}
Deja abierta esa terminal y entra a \url{http://localhost:8000/}. La dirección raíz abre la portada de la versión actual. Si el puerto \texttt{8000} está ocupado, usa otro número tanto en el comando como en la dirección. Para detener el servidor, presiona \texttt{Ctrl+C} en la terminal.

\section{Navegación}
\pasos
\item En la portada, selecciona \boton{Búsquedas} o \boton{Grafos}. Cada catálogo muestra sus métodos agrupados y visibles al cargar.
\item Selecciona un método. La página abre la pestaña correspondiente. Las pestañas de la parte superior del módulo permiten cambiar de método dentro de la misma familia.
\item En escritorio, los menús superiores despliegan accesos directos. En pantalla estrecha, los enlaces a los catálogos siguen visibles y \boton{Menú} abre el panel lateral. \boton{Inicio} y el logo regresan a la portada. \boton{Volver atrás} lleva al nivel anterior.
\item Si necesitas una referencia mientras trabajas, abre \boton{Documentación} para consultar los manuales de la aplicación.
\finpasos
\captura{inicio}{Portada del laboratorio y accesos a los dos catálogos.}{inicio}
\captura{catalogo}{Los métodos de Búsquedas se muestran desde el catálogo.}{catalogo}

\section{Controles y flujo de trabajo}
\begin{tabularx}{\linewidth}{@{}p{.28\linewidth}X@{}}
\toprule
Control & Acción \\
\midrule
\boton{Ejemplo} & Sustituye los datos del método, carga un caso de práctica y ejecuta la operación. Espera si hay una animación. \\
\boton{Crear estructura} & Inicializa la capacidad y el formato de claves. Lee las indicaciones del formulario antes de repetir la operación con datos de trabajo. \\
\boton{Insertar} / \boton{Agregar} & Registra una clave, vértice o arista con los datos escritos. \\
\boton{Buscar} & Recorre la estructura para localizar la clave y muestra la coincidencia o el aviso correspondiente. \\
\boton{Calcular} / \boton{Operar} & Aplica la opción seleccionada a los datos actuales. En algunos métodos se llama \boton{Aplicar algoritmo}. \\
\boton{Eliminar} / \boton{Limpiar} & Elimina el elemento indicado o reinicia la estructura. Revisa la confirmación si aparece. \\
\boton{Guardar} / \boton{Cargar} & Descarga un ejercicio en JSON o recupera uno desde un archivo local. \\
\bottomrule
\end{tabularx}

\chapter{Búsquedas}
\section{Búsquedas internas}
Antes de insertar, crea una estructura y respeta la cantidad de dígitos configurada. Las claves con cero inicial deben escribirse completas.

\metodo{Secuencial}{Define capacidad y longitud, crea la estructura, inserta claves y escribe la que deseas localizar. Pulsa \boton{Buscar}.}{Observa las comparaciones en orden, la posición final y la cantidad de pasos.}
\metodo{Binaria}{Usa la estructura compartida con Secuencial o crea una propia. Escribe la clave y pulsa \boton{Buscar}.}{La vista ordenada muestra cómo se reduce el intervalo de búsqueda después de cada comparación.}
\metodo{Hash}{Configura tamaño y dígitos. Selecciona una función y un método de colisión, crea la estructura e inserta claves.}{Lee la posición inicial y la ruta de resolución cuando dos claves coinciden en la misma posición.}
\metodo{Por residuos}{Elige árbol digital, trie, árbol múltiple o Huffman. En los tres primeros inserta letras de la A a la Z. Para el múltiple, indica el grado y crea el árbol primero. En Huffman, introduce una palabra con al menos dos letras distintas.}{En los árboles se muestra la ruta de bits o grupos de bits. Huffman muestra las frecuencias y los códigos de las letras.}

\section{Búsquedas externas}
Configura la capacidad y los dígitos de las claves en los métodos que usan bloques. Revisa qué bloques o cubetas recorre la animación.
\metodo{Secuencial por bloques}{Crea la estructura, inserta claves y busca una de ellas.}{Sigue el recorrido por bloques y la posición de la clave si aparece.}
\metodo{Binaria por bloques}{Usa los datos compartidos con Secuencial externa o crea la estructura desde esta pestaña.}{El recorrido acota el bloque candidato y luego la posición dentro de ese bloque.}
\metodo{Hash externo}{Selecciona función y estrategia de colisión. Crea la estructura, inserta y después busca.}{Observa la ubicación en bloques y, según la estrategia, la cadena o el área de desbordamiento.}
\metodo{Dinámica}{Elige modo total o parcial, número par de cubetas, registros por cubeta y dígitos. Crea la estructura e inserta claves.}{Consulta la ocupación, los cambios de la estructura y la cubeta identificada al buscar.}

\section{Búsquedas por índices}
\metodo{Un nivel}{Introduce cantidad de registros, tamaño del bloque, tamaño de registro y tamaño de entrada de índice. Los tres tamaños se expresan en bytes. Selecciona primario o secundario y pulsa \boton{Crear estructura}.}{Examina cuántos registros caben por bloque, cuántos bloques se necesitan y cómo se distribuye el índice.}
\metodo{Multinivel}{Introduce los mismos datos y el tipo de índice. Pulsa \boton{Crear estructura}.}{Lee el número de niveles y los accesos estimados. Contrasta la distribución de registros e índices en el diagrama.}

\chapter{Grafos}
\section{Preparación del ejercicio}
\pasos
\item Agrega cada vértice con un nombre distinto. Para un primer caso, usa nombres cortos como \texttt{A}, \texttt{B} y \texttt{C}.
\item Agrega las aristas indicando dos vértices existentes. Si la pantalla pide nombre, asigna uno distinto a cada arista. Si hay dirección o peso, revisa su valor antes de agregar.
\item Elige el tipo de operación, cuando exista un selector. Pulsa el botón de cálculo y examina la notación, el dibujo y las tablas que se presenten.
\finpasos

\section{Operaciones}
\metodo{Un grafo}{Construye el grafo, elige una operación y completa los parámetros que aparezcan. Por ejemplo, para fusionar vértices, indica los dos nombres.}{Compara el grafo original con el resultado y revisa los vértices y aristas resultantes.}
\metodo{Dos grafos}{Construye G1 y G2 en sus formularios. Selecciona una operación, como unión o intersección, y pulsa \boton{Operar}.}{Compara los elementos compartidos o combinados según la operación elegida.}

\section{Árboles como grafos}
\metodo{Centro o bicentro}{Construye un árbol conexo sin ciclos y pulsa \boton{Hallar centro o bicentro}.}{Identifica el vértice central o el par central marcado en el dibujo.}
\metodo{Árbol de expansión mínima}{Construye un grafo conexo ponderado. Completa el peso de cada arista y pulsa el botón de cálculo.}{Distingue las aristas elegidas y consulta el peso total.}
\metodo{Distancia entre dos árboles}{Construye los dos árboles ponderados y pulsa \boton{Hallar distancia}.}{Lee el valor y el desglose que muestra el simulador para esos dos árboles.}

\section{Algoritmos}
\metodo{Función ordinal}{Introduce un grafo dirigido y pulsa \boton{Aplicar algoritmo}.}{Los números identifican el orden asignado a los vértices. Lee el aviso si hay circuitos.}
\metodo{Bellman}{Construye un grafo dirigido ponderado y aplica el algoritmo.}{Examina las distancias desde el origen que usa el simulador. Un aviso señala un ciclo negativo alcanzable cuando corresponda.}
\metodo{Dijkstra}{Agrega aristas dirigidas de peso positivo y aplica el algoritmo.}{Sigue las etiquetas de distancia y predecesor. El camino destacado conecta el primer y último vértice del orden calculado, si existe.}
\metodo{Floyd}{Introduce un grafo dirigido ponderado y aplica el algoritmo.}{Compara la matriz inicial y las iteraciones. Una entrada infinita distingue un par sin camino. Un valor \(-\infty\) señala que un ciclo negativo permite reducir el costo sin límite.}

\section{Matrices}
\metodo{Distancia y circuitos}{Elige distancias, circuitos, circuitos fundamentales, cortes o cortes fundamentales. Construye el grafo con la dirección y pesos que correspondan al ejercicio.}{Lee las filas y columnas según el modo activo. Las matrices de distancias relacionan pares de vértices y las de circuitos o cortes describen conjuntos de aristas.}
\metodo{Incidencia y adyacencia}{Construye un grafo sin peso, selecciona el tipo de matriz y pulsa \boton{Calcular}.}{En incidencia, las filas representan vértices y las columnas aristas. En adyacencia, ambos ejes representan vértices.}

\section{Conjuntos}
\metodo{Independencia}{Construye el grafo y pulsa \boton{Calcular}.}{Examina conjuntos de vértices no adyacentes y distingue los de mayor tamaño.}
\metodo{Dominancia}{Construye el grafo y calcula.}{Lee los conjuntos que cubren a todos los vértices y el número de dominación. Se muestran también las variantes independiente y conexa cuando existen.}
\metodo{Pareamiento}{Agrega vértices y aristas y calcula.}{Identifica los conjuntos de aristas que no comparten extremos y los pareamientos de mayor tamaño.}

\section{Coloreado}
\metodo{Vértices}{Abre la pestaña \boton{Vértices}, construye el grafo y pulsa \boton{Calcular}.}{Dos vértices adyacentes tienen colores distintos. Consulta la cantidad de colores indicada.}
\metodo{Aristas}{Abre \boton{Aristas}, construye el grafo y calcula.}{Las aristas que comparten un extremo tienen colores distintos. Revisa el índice cromático mostrado.}

\chapter{Prácticas ilustradas}
Las capturas de este capítulo se tomaron en el simulador al cargar sus ejemplos integrados. Las instrucciones indican qué datos y resultados deben localizarse en cada vista.

\section{Búsqueda secuencial y comparación con binaria}
\pasos
\item Abre \boton{Búsquedas}, \boton{Búsquedas internas} y \boton{Secuencial}.
\item Pulsa \boton{Ejemplo}. La capacidad se fija en 8, las claves tienen 2 dígitos y la consulta es \texttt{45}.
\item Espera el final de la animación. La fila destacada y el mensaje muestran la posición \textbf{4} y \textbf{4 pasos} (figura~\ref{fig:secuencial}).
\item Cambia a \boton{Búsqueda Binaria}. Escribe \texttt{45} en su campo y pulsa \boton{Buscar} para comparar los recorridos con los mismos datos.
\finpasos
\captura{secuencial}{Ejemplo de búsqueda secuencial con la clave 45 encontrada.}{secuencial}

\clearpage
\section{Búsqueda binaria por bloques}
\pasos
\item En \boton{Búsquedas externas}, abre \boton{Binaria} y pulsa \boton{Ejemplo}.
\item Observa la capacidad de 16 registros, las claves de 3 dígitos y la consulta \texttt{179}.
\item Sigue los bloques inspeccionados y localiza \texttt{179} en el resultado. Compara esta posición con el orden de las claves en la figura~\ref{fig:bloques}.
\finpasos
\captura{bloques}{Búsqueda binaria externa con bloques y resultado de la consulta.}{bloques}

\section{Índice multinivel}
\pasos
\item En \boton{Búsquedas por índices}, abre \boton{Multinivel} y pulsa \boton{Ejemplo}.
\item Lee los parámetros del formulario: 320 registros, bloques de 256 bytes, registros de 32 bytes y entradas de índice de 16 bytes.
\item Con el índice primario se obtienen 8 registros por bloque de datos y 40 bloques. La vista muestra dos niveles de índice y tres accesos contando el bloque de datos.
\finpasos
\captura{indices}{Cálculo de bloques, niveles y accesos del índice primario.}{indices}

\clearpage
\section{Incidencia y adyacencia con el mismo grafo}
\pasos
\item Abre \boton{Grafos}, \boton{Matrices}, y luego \boton{Incidencia y adyacencia}.
\item Elige \boton{Matriz de incidencia} y pulsa \boton{Ejemplo}. Verás tres vértices (A, B y C) unidos por tres aristas. Busca sus nombres en los encabezados de la matriz.
\item Elige \boton{Matriz de adyacencia} y pulsa \boton{Calcular}. Comprueba que ahora los encabezados horizontales también sean vértices (figura~\ref{fig:adyacencia}).
\finpasos
\captura{grafo-matrices}{Grafo dirigido del ejemplo de matrices.}{grafo-matrices}
\captura{incidencia}{Matriz de incidencia del ejemplo integrado.}{incidencia}
\captura{adyacencia}{Matriz de adyacencia calculada con los mismos datos.}{adyacencia}

\clearpage
\section{Caminos y coloreado}
\pasos
\item Abre \boton{Grafos}, \boton{Algoritmos}, \boton{Dijkstra}. Pulsa \boton{Ejemplo}.
\item Compara el camino directo de A a C, de costo 7, con el camino A--B--C, de costo \(2+3=5\). Identifica el segundo en el grafo de resultado (figura~\ref{fig:dijkstra}).
\item Abre \boton{Grafos}, \boton{Coloreado}, \boton{Vértices} y pulsa \boton{Ejemplo}. En el triángulo, cada vértice es adyacente a los otros dos. Verifica que se utilizan tres colores (figura~\ref{fig:coloreado}).
\finpasos
\captura{dijkstra}{Grafo de salida del ejemplo de Dijkstra.}{dijkstra}
\captura{coloreado}{Tres colores en los vértices del triángulo.}{coloreado}

\chapter{Conservar y recuperar ejercicios}
\section{Guardar y volver a cargar}
\pasos
\item Pulsa \boton{Guardar} en el módulo donde tienes el ejercicio. Ubica el archivo JSON en la carpeta de descargas o el destino configurado.
\item Para recuperarlo, abre el mismo módulo y pulsa \boton{Cargar}. Selecciona el archivo que guardaste.
\item Comprueba las claves, vértices o parámetros visibles. Si la pantalla indica que el cálculo está pendiente, vuelve a pulsar \boton{Calcular}, \boton{Operar} o \boton{Aplicar algoritmo}.
\finpasos

\section{Reemplazar los datos}
Guarda los ejercicios que quieras conservar antes de usar \boton{Ejemplo}, limpiar la estructura o recargar la página. Si el objetivo es quitar un solo elemento, indica su clave o nombre y usa \boton{Eliminar}. En un grafo, revisa también las aristas asociadas al vértice eliminado.

\chapter{Solución de problemas}
\begin{tabularx}{\linewidth}{@{}p{.32\linewidth}X@{}}
\toprule
Situación & Qué hacer \\
\midrule
La página local no abre & Comprueba que el servidor esté en ejecución, que lo hayas iniciado desde la raíz del simulador y que el puerto coincida con el de la dirección. \\
No se puede insertar una clave & Revisa que exista una estructura y que queden posiciones libres. Comprueba la cantidad exacta de dígitos y si la clave ya está registrada. \\
Una arista es rechazada & Agrega primero ambos vértices, escribe sus nombres tal como aparecen y completa nombre, dirección o peso cuando se soliciten. \\
No aparece el resultado & Lee el mensaje del método. Algunos algoritmos requieren un grafo conexo, un árbol sin ciclos o pesos positivos. Espera a que termine la animación si está en curso. \\
No se puede cargar un archivo & Selecciona el JSON exportado por el mismo método. Comprueba que el archivo esté completo y que conserve el formato de sus datos. \\
\bottomrule
\end{tabularx}

\chapter{Glosario}
\begin{description}
\item[Clave] Valor que identifica un registro y se utiliza para buscarlo.
\item[Bloque] Grupo de registros tratado como unidad en búsquedas externas o índices.
\item[Cubeta] Espacio donde una función de dispersión ubica uno o varios registros.
\item[Colisión] Coincidencia de dos claves en la misma posición hash.
\item[Índice] Estructura auxiliar que permite localizar registros o bloques.
\item[Vértice y arista] Un vértice es un nodo. Una arista conecta dos nodos y puede tener dirección o peso.
\item[Peso] Valor de una arista que representa un costo o distancia.
\item[Camino y ciclo] Un camino recorre aristas entre vértices. Un ciclo regresa a su vértice inicial.
\item[Árbol] Grafo no dirigido, conexo y sin ciclos.
\item[JSON] Formato de archivo de texto utilizado para intercambiar ejercicios.
\end{description}
\end{document}
